\documentclass{article}
\input{preamble.pre}
\begin{document}

% A box whose last item is a paragraph ending in display math has depth 0 in TeX (the display
% is followed by \belowdisplayskip glue), so its reference point is its very bottom.
\block{VBox ending in display, in H}{
    \showh{\showv{\hsize=200pt\noindent\lorem $$a+b=c$$}}
}
\block{VBox ending in display, in paragraph}{
    \noindent\showpar{\showv{\hsize=200pt\noindent\lorem $$a+b=c$$}}
}
\block{VBox with text after the display, in H}{
    \showh{\showv{\hsize=200pt\noindent\lorem $$a+b=c$$ \lorem}}
}
\block{Bare vbox ending in display (no brace row), in H}{
    \showh{\vbox{\hsize=200pt\noindent\lorem $$a+b=c$$}}
}

% A vtop's reference point is its first item's baseline; when the first item is display math
% (\noindent$$ makes no empty line), the first item is \abovedisplayskip glue, so ht = 0.
\block{VTop starting with display, in H}{
    \showh{\vtop{\hsize=200pt\noindent$$a+b=c$$ \lorem}}
}
\block{VTop starting with text then display, in H}{
    \showh{\vtop{\hsize=200pt\noindent\lorem $$a+b=c$$}}
}

\end{document}
